\documentclass{tietreport}

\usepackage{graphicx}
\usepackage[date=year,%
backend=biber,%
style=alphabetic,%
maxnames=5,%
minnames=3,%
maxalphanames=4,%
minalphanames=3,%
backref=true,%
doi=false,%
isbn=false,%
url=false,%
eprint=false]{biblatex}
\addbibresource{references.bib}

\TitlePageHeader{%
  UCS 503 Software Engineering Project Report%
}

\ProjectTitle{%
  Proposal to Final%
}

\ProjectSubTitle{%
  How to Write a Project Report%
}

\author{%
  \texttt{102XXXXXXX} {Name Title} \\
  \texttt{102XXXXXXX} {Name Title} \\
  \texttt{102XXXXXXX} {Name Title} \\
  \texttt{102XXXXXXX} {Name Title}
}

\TitlePageSubText{%
  Group: \texttt{3XXX} BE Third Year, CSE%
}

\AdvisorName{%
  Dr. NAME OF THE ADVISOR%
}

\TitlePageFooterText{{%
    \large Mid-Semester Evaluation \\[0.2cm]
    \bfseries Computer Science and Engineering
    Department%
  } \\[0.15cm]
  Thapar Institute of Engineering and
  Technology, Patiala%
}

\begin{document}
\maketitle
\tableofcontents

\chapter{Proposal to Project Report}

\section{Overview}

Transforming your proposal into a final project report
is an incremental process of replacing
\textbf{intentions} with \textbf{results}. Since a
proposal serves as a roadmap, the final report becomes
the record of the journey taken.

\subsection{Expanding the Core Structure}

To move from proposal to report, you must shift your
perspective from what you \textit{will} do to what you
\textit{have} done.

\begin{itemize}
\item \textbf{Introduction \& Objectives:} Keep these
  sections but update them to reflect the final state
  of the project. If your objectives shifted during
  development, explain why.
\item \textbf{Methodology to Implementation:} In a
  proposal, the methodology is a plan. In a report,
  this expands into a detailed \textbf{Implementation}
  section. For STEM projects, this includes your final
  system architecture, code modules, and integration
  workflows.
\item \textbf{Results \& Analysis:} This is an entirely
  new section. Replace your ``Deliverables'' list from
  the proposal with actual raw data, prototypes, and
  analyzed findings.
\item \textbf{Conclusion \& Future Work:} Summarize the
  impact of your results and suggest how others might
  improve upon your work.
\end{itemize}

\subsection{Transitioning STEM Technical Content}

For technical projects like AR or 3D modeling, your
methodology evolves into technical documentation.

\begin{itemize}
\item \textbf{From Plan to Architecture:} Move from
  specifying intended tools (e.g., ARCore, Unity) to
  documenting the final \textbf{System Architecture}
  used.
\item \textbf{Documenting Development:} Expand on how
  you handled surface detection, asset management, and
  UI design. Replace ``proposed'' scaling methods with
  documentation of your actual 1:1 metric scaling
  implementation.
\item \textbf{Validation:} Your proposal lists testing
  goals; your report must provide the actual results of
  \textbf{Unit Testing}, \textbf{Performance
    Benchmarking} (like actual FPS achieved), and
  \textbf{User Acceptance Testing} (UAT).
\end{itemize}

\subsection{Incremental Development Steps}

\begin{center}
\begin{tabular}{|l|l|l|}
  \hline
  \textbf{Phase} & \textbf{Proposal} &
  \textbf{Evolution} \\
  \hline
  Initial & SMART Objectives & Validated
  goals; noted deviations. \\
  \hline
  Mid-Project & Methodology \& Tools & Detailed
  System Design and code structure. \\
  \hline
  Execution & Proposed Timeline & Actual log
  of milestones reached. \\
  \hline
  Final & Anticipated Risks & Evaluation \&
  Testing results. \\
  \hline
\end{tabular}
\end{center}

\subsection{Pro-Tips for the Final Report}

\begin{itemize}
\item \textbf{Density Over Fluff:} Just as in the
  proposal, professors value clear, concise information
  over high word counts.
\item \textbf{Visual Documentation:} Replace the
  proposal's flowchart or table with actual screenshots
  of your software, photos of your setup, or graphs of
  your data.
\item \textbf{Addressing Constraints:} Revisit the
  constraints you identified (like lighting or safety)
  and discuss how you successfully mitigated them in
  the final build.
\end{itemize}

\section{More Nuanced Perspective}

Transforming an academic proposal into a final project
report is an incremental process of replacing planned
intentions with documented results. Based on the
provided guide, here is a holistic overview of how to
modify and improvise upon your document through to
completion:

\subsection{Shift from Intent to Execution}

The most significant change is moving from a
``roadmap'' of what you intend to do to a record of
what was actually achieved.

\begin{itemize}
\item \textbf{Methodology to System Architecture:} In a
  STEM project, your methodology evolves from a plan
  into a detailed description of your final
  \textbf{System Architecture} and \textbf{Tech Stack}.
\item \textbf{Module Development:} For the final
  report, expand on how you implemented specific
  features like \textbf{Surface Detection},
  \textbf{Asset Management}, and \textbf{UI Design}.
\item \textbf{Timeline to Project History:} Convert
  your weekly task breakdown into a summary of project
  milestones, noting any shifts in scope or timing.
\end{itemize}

\subsection{Add Validated Results and Testing}

A proposal only anticipates outcomes, whereas a report
must validate them through rigorous testing.

\begin{itemize}
\item \textbf{Performance Benchmarking:} Include
  technical data such as \textbf{Frames Per Second
    (FPS)} and \textbf{latency} measurements to prove
  system stability.
\item \textbf{User Acceptance Testing (UAT):} Document
  the results of ``drift'' analysis---checking if
  virtual objects remain stable in their environment
  over time.
\item \textbf{Unit Testing:} Provide a log of
  individual feature tests, such as verifying if
  rotation or placement logic worked as intended.
\end{itemize}

\subsection{Refine and Expand Core Sections}

\begin{itemize}
\item \textbf{Introduction/Problem Statement:} Refine
  this section (ideally 250--400 words) to ensure it
  accurately reflects the final problem addressed by
  your completed work.
\item \textbf{Objectives:} Revisit your \textbf{SMART
    objectives}. In your report, state clearly whether
  each objective was met, partially met, or modified.
\item \textbf{Resources:} Update your resource list to
  reflect what was actually used, such as specific
  \textbf{AR libraries} (ARCore, ARKit) or hardware
  like \textbf{LiDAR-enabled devices}.
\end{itemize}

\subsection{Address Real-World Constraints}

While the proposal mentions anticipated risks,
the final report should discuss how you managed
them during development.

\begin{itemize}
\item \textbf{Environmental Constraints:} Discuss how
  the final system performed in low-light conditions or
  on reflective surfaces.
\item \textbf{Technical Challenges:} Detail how you
  solved issues like \textbf{coordinate mapping} (1:1
  scaling) or \textbf{physics engine} clipping.
\end{itemize}

\subsection{Final Formatting Adjustments}

\begin{itemize}
\item \textbf{Tone:} Maintain a professional tone using
  active verbs to describe your completed actions.
\item \textbf{Density:} Ensure the final document
  remains dense with information rather than ``fluff'',
  focusing on proving the project's academic or
  scientific soundness.
\item \textbf{Visuals:} Replace the proposal's
  placeholder tables with final data visualizations,
  flowcharts of your actual technical workflow, and
  screenshots of the finished product.
\end{itemize}

\section{Example Details}

These are the ``meat'' of your final report and
transition your proposal's ``Methodology'' into a
documented reality.

\subsection{Section: Technical System
  Architecture}

In the proposal, you listed tools. In the report, you
explain how they interact. You should replace your
methodology flowchart with a high-level system diagram.

\textbf{Drafting Template:}

\begin{itemize}
\item \textbf{The Tech Stack:} ``The project was
  developed using \textbf{Unity 2022.3} as the primary
  engine, utilizing the \textbf{ARCore SDK} for spatial
  mapping.  C\# was used for logic scripting, while
  \textbf{Blender} served as the pipeline for 1:1
  metric-scale 3D modeling.''
\item \textbf{The Workflow:} Describe the data
  flow. ``The system initiates by polling the device
  camera for feature points. Once a horizontal plane is
  detected via \textbf{Raycasting}, the coordinate
  system is anchored to a 0,0,0 origin point to ensure
  stability.''
\end{itemize}

\subsection{Section: Module-by-Module Implementation}

Break down the ``doing'' part of your project into
logical modules. This proves the complexity of your
work.

\begin{itemize}
\item \textbf{Module A: Environment Scanning \& Plane
    Detection}
  \begin{itemize}
  \item \textit{Details:} How did you handle lighting?
    Did you use Depth APIs?
  \item \textit{Improvement:} ``While the proposal
    suggested basic plane detection, the implementation
    utilized \textbf{Light Estimation APIs} to
    dynamically adjust the shaders on virtual objects,
    increasing realism.''
  \end{itemize}
\item \textbf{Module B: Interaction Logic}
  \begin{itemize}
  \item \textit{Details:} How does the user move or
    rotate objects?
  \item \textit{Improvement:} ``We implemented a
    \textbf{Gesture Recognition} module that translates
    screen-space touch inputs into world-space
    transformations, restricted by collision boxes to
    prevent `ghosting' through walls.''
  \end{itemize}
\end{itemize}

\subsection{Section: Testing and Validation}

This is the most critical addition for a final
report. Use the SMART objectives from your proposal to
create a ``Pass/Fail'' or ``Result'' table.

\textbf{Key Metrics to Include:}

\begin{itemize}
\item \textbf{Performance Benchmarking:} ``The
  application maintained a consistent \textbf{55--60
    FPS} on a Samsung S22, though performance dropped
  to 40 FPS in low-light environments due to increased
  CPU load from feature-point searching.''
\item \textbf{Drift Analysis:} ``In a 10-minute stress
  test, the virtual object's anchor point drifted by
  less than 2cm, meeting our accuracy objective.''
\end{itemize}

\chapter{Examples: \LaTeX{} Capabilities}

\section{Tables and Structured Data}

\LaTeX{} excels at creating professional tables with
precise alignment and
formatting. Table~\ref{tab:performance}~\cite{latex2e}
shows typical performance metrics captured during
testing phases.

\begin{table}[h]
\centering
\label{tab:performance}
\begin{tabular}{|l|r|r|r|}
  \hline
  \textbf{Component} & \textbf{Load (ms)} &
  \textbf{FPS} & \textbf{Memory (MB)} \\
  \hline
  Rendering Engine & 2.5 & 60 & 128 \\
  \hline
  Physics Simulation & 3.1 & 55 & 64 \\
  \hline
  Asset Management & 1.8 & 60 & 96 \\
  \hline
  UI Framework & 0.9 & 60 & 32 \\
  \hline
\end{tabular}

\caption{Performance Metrics}
\end{table}

\section{Figures and Visual Content}

Figures can be embedded using the graphicx package.
Figure~\ref{fig:sample} represents a typical workflow
diagram (placeholder).

\begin{figure}[h]
  \centering%
  \includegraphics[width=4in]{%
    sample-report-target-p1.png}
  \caption{Sample title page layout%
    \label{fig:sample}}
\end{figure}

Multiple figures can be arranged side-by-side or
stacked vertically using minipage environments,
enabling complex document layouts.

\section{Citations and References}

Academic reports benefit from proper citation
management.  This section demonstrates how to cite
sources using \LaTeX{}'s bibliography system.
According to the documentation~\cite{lamport94},
\LaTeX{} provides powerful tools for document
preparation.

The combination of structured markup and automated
formatting ensures consistency across all citations and
references throughout the document~\cite{knuth1984}.

\subsection{Using Biblatex for Advanced Citation
  Management}

The \texttt{biblatex} package provides modern,
flexible citation management with powerful
formatting options. Biblatex offers multiple
citation styles and commands:

\begin{itemize}
\item \texttt{\textbackslash cite\{\}}:
  Parenthetical citations: ``This is well
  documented \cite{lamport94,knuth1984}.''
\item \texttt{\textbackslash textcite\{\}}:
  Textual citations: ``According to
  \textcite{latex2e}, LaTeX2e provides robust
  tools for professional documents.''
\item \texttt{\textbackslash parencite\{\}}:
  Alternative parenthetical style.
\item \texttt{\textbackslash footcite\{\}}:
  Footnote citations.
\end{itemize}

Biblatex requires the \texttt{biber} backend,
providing superior bibliography management
compared to older \texttt{natbib} package.

\subsection{Integrating Zotero with Biblatex}

Zotero is a free, open-source reference
manager that integrates seamlessly with
biblatex via BetterBibTeX. The workflow
is straightforward:

\begin{enumerate}
\item Install Zotero and the BetterBibTeX
  extension from \texttt{retorque.re}.
\item Create or import your reference
  library into Zotero.
\item Export your library as a
  \texttt{.bib} file (BibTeX format).
\item Reference the file in your LaTeX
  preamble using
  \texttt{\textbackslash addbibresource\{%
  references.bib\}}.
\item Cite sources using
  \texttt{\textbackslash cite\{\}},
  \texttt{\textbackslash textcite\{\}},
  or \texttt{\textbackslash footcite\{\}}.
\item Compile with: \texttt{pdflatex}
  $\rightarrow$ \texttt{biber}
  $\rightarrow$ \texttt{pdflatex}.
\end{enumerate}

Zotero automatically generates citation
keys. BetterBibTeX auto-exports your
library to \texttt{.bib} file, keeping
references synchronized with your Zotero
collection automatically.

\chapter{Conclusion}

To make this specific to your actual project, start by
looking into:

\begin{enumerate}
\item \textbf{What is the specific title or goal of
    your project?} (e.g., An AR furniture app? A data
  analysis tool?)
\item \textbf{What was the hardest technical hurdle you
    faced?} (We’ll together write the ``Challenges \&
  Mitigations'' section for this).
\end{enumerate}

\printbibliography[heading=bibintoc]

\end{document}
